\chapter{Control Interfaces}

\section*{학습 목표}
이 장은 VLA policy가 출력한 $\act_t$가 어떤 제어 interface를 거쳐 실제 robot command $\ctrlcmd_t$가 되는지 다룬다. Chapter 2가 state와 observation의 차이를 고정했다면, Chapter 3은 action proposal과 controller command의 차이를 고정한다.

\begin{enumerate}
  \item VLA action representation과 controller interface를 분리한다.
  \item joint, Cartesian, waypoint, trajectory, impedance, MPC reference의 차이를 설명한다.
  \item $\act_t \rightarrow \ctrlcmd_t$ mapping이 safety와 evaluation에 미치는 영향을 정식화한다.
  \item Chapter 13의 action representation 논의를 위한 control-side 기준을 제공한다.
\end{enumerate}

\section{Control interface가 필요한 이유}

VLA policy가 ``action을 출력한다''는 말은 충분하지 않다. Action이 무엇인지, 어떤 rate로 나오며, 어떤 coordinate frame을 쓰고, 어떤 controller가 그 action을 해석하는지 명시해야 한다. 실제 로봇은 neural network output을 그대로 실행하지 않는다. 대부분의 경우 VLA output은 controller reference, waypoint, target pose, short trajectory, 또는 high-level skill call이다 \cite{khatib1987,hogan1985,mayne2000}.

Chapter 1의 notation을 다시 쓰면, VLA policy는
\begin{equation}
  \act_t = \policy(\obs_t, \conf_t, \task, \bel_t)
  \label{eq:ch3-policy}
\end{equation}
를 출력한다. 하지만 robot driver는 보통 $\act_t$가 아니라
\begin{equation}
  \ctrlcmd_t = \controller(\act_t,\hat{\state}_t,\mathcal{C}_t)
  \label{eq:ch3-controller-map}
\end{equation}
를 받는다. 여기서 $\controller$는 단순한 software wrapper가 아니다. 그것은 action representation을 물리적 command로 해석하는 interface contract이다.

\begin{definitionbox}{Control interface}
Control interface는 policy-level action $\act_t$가 robot hardware에 전달되기 전에 통과하는 해석 경계이다. 이 경계는 coordinate frame, unit, rate, limit, tracking law, safety filter, failure semantics를 포함한다.
\end{definitionbox}

\section{Joint-space interfaces}

가장 직접적인 interface는 joint-space command이다. Robot configuration을 $\conf\in\R^n$이라 할 때, VLA 또는 downstream module은 target joint position, joint velocity, 또는 torque를 낼 수 있다.

\subsection{Joint position target}

Joint position interface는 다음과 같이 쓸 수 있다.
\begin{equation}
  \act_t = \conf^{ref}_{t+H},
  \qquad
  \ctrlcmd_t = K_p(\conf^{ref}_{t+H}-\conf_t)-K_d\dot{\conf}_t.
  \label{eq:joint-position-control}
\end{equation}

이 interface는 단순하고 많은 robot API와 잘 맞는다. 그러나 VLA가 joint target을 직접 내면 morphology dependency가 매우 커진다. 같은 task라도 robot arm의 joint dimension, joint limit, kinematic chain이 다르면 action space가 달라진다. Multi-embodiment VLA에서 joint-space output이 어려운 이유가 여기에 있다.

\subsection{Joint velocity target}

Velocity command는 다음처럼 표현된다.
\begin{equation}
  \act_t = \dot{\conf}^{ref}_t,
  \qquad
  \ctrlcmd_t = \dot{\conf}^{ref}_t.
  \label{eq:joint-velocity-control}
\end{equation}
단순해 보이지만, velocity command는 integration drift와 limit handling이 중요하다. Velocity를 오래 적분하면 joint limit, self-collision, workspace boundary를 넘을 수 있다. 따라서 velocity interface는 safety projection과 함께 쓰여야 한다.

\subsection{Torque command}

Torque-level interface는 가장 낮은 수준의 command이다.
\begin{equation}
  \act_t = \bm{\tau}^{ref}_t,
  \qquad
  \ctrlcmd_t = \bm{\tau}^{ref}_t.
  \label{eq:torque-interface}
\end{equation}
Torque command는 표현력이 크지만 risk도 크다. 고주파 안정성, latency, actuator saturation, contact instability를 직접 다루어야 한다. VLA policy가 torque를 직접 출력하려면 control frequency, safety certification, dynamics generalization 문제가 매우 어려워진다.

\begin{warningbox}{Torque output의 위험}
Torque는 robot hardware에 가장 가까운 action이다. 학습 policy가 torque를 직접 내면 controller가 흡수하던 안정성 책임이 policy로 이동한다. 이 경우 inference delay, out-of-distribution state, contact uncertainty가 곧바로 safety issue가 된다.
\end{warningbox}

\section{Cartesian interfaces}

대부분의 manipulation task는 joint보다 task-space에서 표현하기 쉽다. End-effector pose를 $\bm{x}_{ee}\in SE(3)$ 또는 local coordinate vector로 쓰면, VLA output은 target pose 또는 delta pose가 될 수 있다.

\subsection{End-effector target pose}

End-effector target pose interface는 다음처럼 쓸 수 있다.
\begin{equation}
  \act_t = \bm{x}^{ref}_{ee,t+H}.
  \label{eq:ee-target}
\end{equation}
Controller는 inverse kinematics 또는 operational-space control을 통해 이를 joint command로 바꾼다.
\begin{equation}
  \dot{\conf}^{ref}_t
  = J(\conf_t)^\dagger
  K_x(\bm{x}^{ref}_{ee,t+H}-\bm{x}_{ee,t}).
  \label{eq:ik-command}
\end{equation}
여기서 $J(\conf_t)$는 Jacobian, $J^\dagger$는 pseudo-inverse이다.

이 interface는 robot morphology를 일부 숨길 수 있다. Policy는 end-effector motion을 내고, robot-specific kinematics는 controller가 처리한다. 그러나 singularity, joint limit, collision, orientation ambiguity는 여전히 controller가 해결해야 한다.

\subsection{Delta pose}

Delta pose interface는 VLA에서 자주 사용된다.
\begin{equation}
  \act_t = \Delta \bm{x}_{ee,t},
  \qquad
  \bm{x}^{ref}_{ee,t+1} = \bm{x}_{ee,t}\oplus \Delta \bm{x}_{ee,t}.
  \label{eq:delta-pose}
\end{equation}
여기서 $\oplus$는 pose composition을 의미한다. Delta pose는 local correction에 유리하지만, scale과 coordinate frame에 민감하다. Camera frame 기준 delta인지, base frame 기준 delta인지, end-effector frame 기준 delta인지 명시하지 않으면 같은 vector가 완전히 다른 motion이 된다.

\begin{table}[ht]
\centering
\caption{Cartesian interface에서 반드시 명시해야 할 항목}
\label{tab:cartesian-contract}
\small
\begin{tabularx}{\textwidth}{p{32mm}X}
\toprule
항목 & 이유 \\
\midrule
Frame & base, camera, end-effector frame 중 무엇인지에 따라 motion 방향이 달라짐 \\
Unit and scale & normalized action인지 metric meter/radian인지 구분 필요 \\
Rate & delta가 control step당 변화인지 VLA inference step당 변화인지 명시 필요 \\
Orientation representation & Euler, quaternion, axis-angle, 6D representation의 singularity와 normalization 문제 \\
Limit handling & workspace, velocity, acceleration, joint limit projection 필요 \\
\bottomrule
\end{tabularx}
\end{table}

\section{Waypoint and trajectory interfaces}

Waypoint interface는 policy가 단일 command 대신 중간 목표를 제안하는 구조이다.
\begin{equation}
  \act_t = \bm{w}_{t:t+K}
  =
  \{\bm{x}^{ref}_{1},\ldots,\bm{x}^{ref}_{K}\}.
  \label{eq:waypoint-sequence}
\end{equation}
Motion planner 또는 trajectory optimizer는 waypoint를 연결해 시간 매개화된 trajectory를 만든다.
\begin{equation}
  \xi^\star
  =
  \arg\min_{\xi}
  J(\xi;\bm{w}_{t:t+K})
  \quad
  \text{s.t.}\quad
  \xi(\tau)\in\mathcal{X}_{free},\;
  \dot{\xi}(\tau)\in\mathcal{V},\;
  \xi(0)=\state_t.
  \label{eq:trajectory-optimization}
\end{equation}

Waypoint interface의 장점은 policy와 motion planning의 역할을 분리한다는 점이다. Policy는 semantic target 또는 local waypoint를 내고, planner는 collision-free path와 timing을 만든다. 단점은 planner failure가 policy failure처럼 보일 수 있다는 점이다. Evaluation log는 waypoint feasibility, planner failure, controller tracking failure를 분리해야 한다.

\section{Impedance and operational-space interfaces}

Contact-rich manipulation에서는 position tracking만으로 부족하다. 로봇이 물체를 밀거나, 잡거나, 삽입하거나, 사람 가까이에서 움직일 때는 compliance가 필요하다.

\subsection{Impedance target}

Impedance control은 end-effector가 가상의 spring-damper처럼 동작하도록 한다.
\begin{equation}
  \bm{F}_{cmd}
  =
  K_x(\bm{x}^{ref}-\bm{x})
  +D_x(\dot{\bm{x}}^{ref}-\dot{\bm{x}})
  +\bm{F}^{ff}.
  \label{eq:impedance}
\end{equation}
이때 VLA policy는 position target뿐 아니라 stiffness, damping, force reference를 제안할 수도 있다.
\begin{equation}
  \act_t = (\bm{x}^{ref}_t, K_{x,t}, D_{x,t}, \bm{F}^{ff}_t).
  \label{eq:impedance-action}
\end{equation}

이 interface는 manipulation에 강력하지만 safety risk도 크다. Stiffness를 잘못 선택하면 contact force가 커지고, damping이 부족하면 oscillation이 생긴다. 따라서 impedance parameter를 VLA가 직접 내는 경우에는 safety bound가 필요하다.

\subsection{Operational-space control}

Operational-space control은 task-space force를 joint torque로 변환한다.
\begin{equation}
  \bm{\tau}
  =
  J(\conf)^\top \bm{F}_{cmd}
  + \bm{N}(\conf)^\top \bm{\tau}_{null}
  + \bm{g}(\conf).
  \label{eq:osc}
\end{equation}
여기서 $\bm{N}$은 null-space projection이다. 이 구조는 VLA policy가 task-space target을 내고, controller가 robot-specific dynamics를 처리하는 좋은 예이다.

\section{MPC reference interface}

Model predictive control(MPC)은 VLA output을 reference 또는 constraint로 받아 finite horizon optimization을 수행할 수 있다.

\begin{equation}
\begin{aligned}
  \ctrlcmd_{t:t+H-1}^\star
  =
  \arg\min_{\ctrlcmd_{t:t+H-1}}
  &\sum_{k=0}^{H}
  \ell_x(\state_{t+k}, \act_t)
  +\sum_{k=0}^{H-1}\ell_u(\ctrlcmd_{t+k}) \\
  \text{s.t.}\quad
  &\state_{t+k+1}=f(\state_{t+k},\ctrlcmd_{t+k}),\\
  &\state_{t+k}\in\mathcal{X}_{safe},\quad
  \ctrlcmd_{t+k}\in\mathcal{U}.
\end{aligned}
\label{eq:mpc-vla}
\end{equation}

여기서 $\act_t$는 target pose, cost parameter, forbidden region, subgoal, 또는 terminal set으로 들어갈 수 있다. MPC interface의 장점은 constraint와 prediction을 명시적으로 다룰 수 있다는 점이다. 단점은 model mismatch, optimization time, solver failure를 관리해야 한다는 점이다.

\section{Safety projection}

VLA output을 그대로 controller에 넣지 않고 safety set으로 project할 수 있다. 가장 단순한 형태는 다음이다.
\begin{equation}
  \tilde{\act}_t
  =
  \arg\min_{\act\in\mathcal{A}_{safe}(\hat{\state}_t)}
  \|\act-\act_t\|^2.
  \label{eq:safety-projection}
\end{equation}
이후 controller는 $\tilde{\act}_t$를 사용한다.
\begin{equation}
  \ctrlcmd_t = \controller(\tilde{\act}_t,\hat{\state}_t,\mathcal{C}_t).
  \label{eq:projected-control}
\end{equation}

이 방식은 action space가 metric structure를 가질 때 유용하다. Continuous delta나 waypoint는 projection이 비교적 자연스럽다. 반면 action token은 물리적 distance metric이 명확하지 않기 때문에 projection이 어렵다. 이것이 action representation과 safety가 연결되는 지점이다.

\section{Control interface별 failure mode}

\begin{table}[ht]
\centering
\caption{Control interface별 대표 failure mode}
\label{tab:control-failure}
\small
\begin{tabularx}{\textwidth}{p{28mm}X X}
\toprule
Interface & 강점 & 대표 failure mode \\
\midrule
Joint position & robot API와 단순 연결 & morphology-specific, joint limit violation \\
Joint velocity & reactive correction 가능 & integration drift, accumulated unsafe motion \\
Torque & full dynamics control 가능 & latency와 OOD state에서 unstable behavior \\
End-effector pose & task-space 표현이 직관적 & IK failure, singularity, frame mismatch \\
Delta pose & local closed-loop correction에 적합 & scale/frame ambiguity, drift \\
Waypoint & planner와 결합 가능 & infeasible waypoint, planner/controller blame ambiguity \\
Impedance target & contact-rich task에 적합 & stiffness/force unsafe setting \\
MPC reference & constraint와 horizon을 명시 가능 & model mismatch, solver latency \\
\bottomrule
\end{tabularx}
\end{table}

이 표는 Chapter 13의 action taxonomy를 평가할 때 그대로 다시 사용된다. 어떤 action representation도 단독으로 좋거나 나쁘지 않다. 중요한 것은 task, controller, safety monitor, data log가 같은 interface contract를 공유하는가이다.

\section{Evaluation log에 남겨야 할 control metadata}

VLA paper가 real robot result를 보고할 때 control interface metadata가 빠지면 결과 해석이 불가능해진다. 최소한 다음 항목은 필요하다.

\begin{itemize}
  \item action type: token, delta pose, waypoint, velocity, torque 등
  \item coordinate frame: base, camera, end-effector, world
  \item VLA inference rate와 controller rate
  \item downstream controller type
  \item action scaling과 normalization
  \item limit handling과 safety projection
  \item planner 또는 controller failure count
  \item human intervention과 reset policy
\end{itemize}

이 metadata가 있어야 benchmark success와 robot capability를 분리할 수 있다. 같은 success rate라도 torque-level policy와 waypoint+MPC system은 전혀 다른 evidence를 제공한다.

\begin{casebox}{Case study: 같은 action, 다른 controller claim}
``move 5 cm right''라는 action은 Cartesian delta controller에서는 local servo command이고, waypoint planner에서는 collision-checked target이며, impedance controller에서는 contact-compliant reference가 된다. 따라서 action string이 같아도 system claim은 다르다. 논문은 action syntax가 아니라 downstream controller를 밝혀야 한다.
\end{casebox}

\chapterwork
{Operational-space control, impedance control, MPC를 Chapter 13의 action representation taxonomy와 함께 읽어라 \cite{khatib1987,hogan1985,mayne2000}.}
{왜 VLA output $\act_t$를 robot command $\ctrlcmd_t$와 동일시하면 safety와 evaluation 해석이 약해지는가?}
{하나의 tabletop task에 대해 delta pose, waypoint, impedance target 세 가지 controller interface를 설계하고 각각 필요한 metadata를 비교하라.}

\section{Roboticist's Takeaway}

\begin{definitionbox}{Chapter 3 Takeaway}
VLA action은 output tensor가 아니라 control interface object이다. $\act_t$의 의미는 downstream controller $\controller$, state estimate $\hat{\state}_t$, safety constraint $\mathcal{C}_t$, command rate가 함께 정의할 때만 결정된다. 따라서 Real VLA 논문을 읽을 때는 model architecture만 보지 말고, $\act_t$가 어떤 controller를 통해 $\ctrlcmd_t$가 되는지 확인해야 한다.
\end{definitionbox}

다음 장은 이 control interface가 contact, force, manipulation과 만날 때 어떤 문제가 생기는지 다룬다.
